\documentclass[11pt]{article}

\usepackage[T1]{fontenc}
\usepackage{lmodern}
\usepackage[margin=0.92in]{geometry}
\usepackage{amsmath,amssymb,amsthm,mathtools}
\usepackage{microtype}
\usepackage{graphicx}
\usepackage[hidelinks]{hyperref}
\usepackage{enumitem}

\setlist{nosep}
\newtheorem{theorem}{Theorem}
\newtheorem{lemma}{Lemma}
\theoremstyle{remark}
\newtheorem{remark}{Remark}

\newcommand{\V}{V}
\newcommand{\one}{\mathbf 1}
\newcommand{\Leb}{\operatorname{Leb}}
\newcommand{\DGF}{\operatorname{DGF}}

\title{Central Area Drifts Obstruct Intrinsic Uniqueness\\of Rectangular Young Surface Lifts}
\author{Independent research note}
\date{August 17, 2026}

\begin{document}
\maketitle

\begin{abstract}
Lee asks whether rectangular increment regularity can yield a unique lift of
a surface in the two-parameter Young regime.  We give a negative answer to
intrinsic uniqueness under the paper's stated rough-surface axioms.  For every
$\rho\in(1/2,2/3]$, the zero surface in $\mathbb R^3$ has a one-parameter
family of distinct full rectangular $\rho$-Hölder multiplicative
double-group lifts sharing all path data and all surface levels through two.
The ambiguity is a degree-three Jacobi cycle multiplied by Euclidean rectangle
area.  It satisfies even the stronger mixed rectangular estimate suggested in
the paper's appendix.  Additional normalization or naturality could still
select the canonical Young-integral lift.
\end{abstract}

\section*{Packet status}

\begin{description}[leftmargin=3.0cm,style=nextline]
\item[Run] \texttt{fa\_banach\_001}, lane 13.
\item[Source] arXiv:2406.16857, Remark~5.23 and the first question on PDF page~56 \cite{Lee}.
\item[Classification] Candidate counterexample, likely valid subject to human review.
\item[Scope] Intrinsic uniqueness among admissible rectangular rough lifts for
$1/2<\rho\le2/3$; canonical selection under extra axioms is not excluded.
\end{description}

\begin{center}
\includegraphics[width=0.97\textwidth]{figures/open_problem_crop.png}
\end{center}

\section{Question and answer}

The source asks whether an analogue of its surface extension theorem can use
rectangular regularity to ``define a unique lift'' of
$X\in C^{\square,\rho}([0,1]^2,\V)$ throughout the Young regime
$\rho\in(1/2,1]$ \cite[Section~6]{Lee}.  Remark~5.23 makes the missing issue
precise: the constructed extension is unique only after level three has been
supplied, and the paper does not prove that level three is determined by the
underlying surface.

We use the source's free crossed module
$(\delta:T_1(\V)\to T_0(\V),\triangleright)$ and its group-like elements
$(G_1,G_0)$.  The product on $T_1$ is denoted by $*$.  For a rectangle
$Q=[s_1,s_2]\times[t_1,t_2]$, write
\[
 |Q|=(s_2-s_1)(t_2-t_1).
\]

\begin{theorem}[central-area nonuniqueness]\label{thm:main}
Let $\V=\mathbb R^3$ and $1/2<\rho\le2/3$.  The constant surface
$X\equiv0$ admits a family $(\mathbf X^\lambda)_{\lambda\in\mathbb R}$ of
full multiplicative double-group functionals in
$\DGF^{\square,\rho}(G_1,G_0)$ such that:
\begin{enumerate}[label=(\roman*)]
\item every horizontal and vertical path component is the identity;
\item every surface component through level two is zero;
\item the family members are pairwise distinct.
\end{enumerate}
Thus the underlying rectangular $\rho$-Hölder surface, even together with
all boundary paths and all data through level two, does not intrinsically
determine its lift under the paper's regularity conditions.
\end{theorem}

\section*{Idea of the proof}

The first nontrivial kernel of the crossed-module boundary contains the
degree-three Jacobi cycle
\[
 \Omega=[e_1,e_2\wedge e_3]+[e_2,e_3\wedge e_1]
       +[e_3,e_1\wedge e_2].
\]
Jacobi gives $\delta\Omega=0$, while the six terms in its free-bimodule
expansion show $\Omega\ne0$.  The Peiffer identities make kernel elements
square-zero against all of $T_1$.  Consequently
$\exp_*(a\Omega)=\one+a\Omega$, and multiplication of these exponentials
just adds $a$.  Assigning $a=\lambda|Q|$ to a rectangle therefore obeys
both horizontal and vertical concatenation because area is additive in both
directions.  Finally, $|Q|$ has exactly the critical two-dimensional scaling:
it meets the degree-three surface bound precisely up to $\rho=2/3$.

\section{The degree-three kernel cycle}

Let $e_1,e_2,e_3$ be the standard orthonormal basis.  The Lie action is
embedded in the free bimodule by
\[
 [x,u\wedge v]=x\cdot(u\wedge v)-(u\wedge v)\cdot x.
\]

\begin{lemma}\label{lem:omega}
The element $\Omega$ above belongs to the degree-three part of Kapranov's
semiabelianized Lie algebra inside $T_1(\V)$, is nonzero, and satisfies
\[
 \delta\Omega=0,\qquad \|\Omega\|=\sqrt6,
 \qquad \Omega*F=F*\Omega=0\quad(F\in T_1(\V)).
\]
In particular, $\exp_*(a\Omega)=\one+a\Omega\in G_1$ for every $a\in\mathbb R$.
\end{lemma}

\begin{proof}
The boundary intertwines the Lie action, hence
\[
 \delta\Omega=[e_1,[e_2,e_3]]+[e_2,[e_3,e_1]]
                 +[e_3,[e_1,e_2]]=0
\]
by Jacobi.  On expanding the three action brackets, one obtains six distinct
signed vectors of the orthonormal basis
\[
 \V\otimes\Lambda^2\V\ \oplus\ \Lambda^2\V\otimes\V.
\]
Every Peiffer generator has total degree at least $2+2=4$, so the Peiffer
quotient imposes no relation in degree three.  Therefore $\Omega\ne0$ and
$\|\Omega\|^2=6$.

The second Peiffer identity gives
$E*F=\delta(E)\cdot F=E\cdot\delta(F)$.  Taking $E=\Omega$ in the first
form and $F=\Omega$ in the second proves both zero-product identities.
Thus every positive power $\Omega^{*m}$ with $m\ge2$ vanishes.  Since
$\Omega$ lies in the Lie algebra used to define $G_1$,
$\exp_*(a\Omega)=\one+a\Omega$ is group-like at every truncation.  Its
finite graded support also places it in the completed group $G_1$.
\end{proof}

\section{The full multiplicative lift}

For $\lambda\in\mathbb R$, give $X\equiv0$ identity-valued horizontal and
vertical path components and define its surface component on every rectangle
$Q$ by
\begin{equation}\label{eq:drift}
 \mathbf X_Q^\lambda
 =\exp_*(\lambda|Q|\Omega)
 =\one+\lambda|Q|\Omega.
\end{equation}
The boundary condition holds because
$\delta(\mathbf X_Q^\lambda)=\exp(\lambda|Q|\delta\Omega)=\one$, which is
the boundary signature of the zero surface.  If $Q_1,Q_2$ are horizontally
adjacent, all path actions in the double-group composition are trivial, and
Lemma~\ref{lem:omega} gives
\[
 \mathbf X_{Q_2}^\lambda*\mathbf X_{Q_1}^\lambda
 =\one+\lambda(|Q_1|+|Q_2|)\Omega
 =\mathbf X_{Q_1\cup Q_2}^\lambda.
\]
The same calculation applies to vertical adjacency.  Hence
\eqref{eq:drift} is a full multiplicative double-group functional, without
appealing to an extension theorem.  Distinct parameters give distinct values
on $[0,1]^2$, while all levels below three agree.

\section{Regularity at the critical two-dimensional scale}

It remains to check the analytic hypothesis.  Put $a=s_2-s_1$, $b=t_2-t_1$
and $\sigma=\rho/2$.  The only nonzero positive homogeneous component of
\eqref{eq:drift} is
\[
 \| (\mathbf X_Q^\lambda)^{(3)}\|
 =|\lambda|\sqrt6\,ab.
\]
The source's level-three surface regularity polynomial contains the positive
$q=3$ term
\[
 \frac{2^{6\sigma}}{\beta}
 \frac{(C_\omega a)^{3\sigma}(C_\omega b)^{3\sigma}}
      {\Gamma(3\sigma+1)^2}.
\]
Because $3\sigma=3\rho/2\le1$ and $a,b\in[0,1]$, we have
$ab\le a^{3\sigma}b^{3\sigma}$.  Choosing the Hölder-control constant
$C_\omega$ large enough absorbs $|\lambda|\sqrt6$ and the fixed coefficient.
All path bounds and all other surface-level bounds are zero.  Therefore
$\mathbf X^\lambda\in\DGF^{\square,\rho}(G_1,G_0)$.

The example also meets the sharper mixed rectangular scaling mentioned in
the appendix:
\begin{equation}\label{eq:strong}
 ab\le a^\rho b^{2\rho}+a^{2\rho}b^\rho.
\end{equation}
Indeed, if $a\le b$, then
\[
 \frac{ab}{a^\rho b^{2\rho}}
 =a^{1-\rho}b^{1-2\rho}
 \le b^{2-3\rho}\le1;
\]
if $b\le a$, use the second term symmetrically.  This proves
Theorem~\ref{thm:main}.

\section{Interpretation and limitations}

On square cells of side $\varepsilon$, the drift has size
$\varepsilon^2$, while degree-three rough regularity has scale
$\varepsilon^{3\rho}$.  The threshold $3\rho=2$ explains both why the
construction works throughout the unresolved interval $(1/2,2/3]$ and why
it does not conflict with the source's existing uniqueness from level two
for $\rho>2/3$.

The conclusion concerns \emph{intrinsic} uniqueness inside the stated
rectangular DGF class.  The source's Young-integral formula remains a valid
and natural canonical choice; the nonzero lifts above do not equal the smooth
surface signature of the zero surface.  Accordingly, the example does not
exclude a theorem that adds smooth-approximation consistency, a locality or
naturality axiom, or a quantitative normalization tying the higher-level
control to the norm of $X$.  It proves that rectangular regularity and the
existing multiplicativity/boundary axioms alone cannot supply uniqueness.

\section{Verification and novelty audit}

The accompanying exact-arithmetic script expands $\Omega$, verifies its zero
boundary and squared norm six, and checks both area additivity laws on
rational grids.  It also tests \eqref{eq:strong} on $16{,}384$ exponent and
aspect-ratio cases; all checks pass.  These computations corroborate, but do
not replace, the symbolic proof.

On August 17, 2026, the run indexes and local arXiv corpus were searched by
the source id/title and combinations of \emph{unique lift}, \emph{rectangular
Hölder}, \emph{central kernel}, and \emph{Jacobi}.  Bounded web searches used
the same terms together with rough- and multiplicative-surface-signature
variants.  No later explicit answer or this construction was located.  This
is a bounded novelty check, not an exhaustive bibliographic claim.

\begin{thebibliography}{9}
\bibitem{Lee}
D. Lee,
\emph{The Surface Signature and Rough Surfaces},
arXiv:2406.16857v2 (2024).
\end{thebibliography}

\end{document}
